% !TeX program = lualatex
\documentclass[a5paper,11pt,openany]{article}
\usepackage{orcidlink}

\usepackage{blindtext}
%\usepackage[style]{abstract}

\usepackage[top=2cm,bottom=2cm,bindingoffset=0cm]{geometry}

\usepackage[svgnames]{xcolor} % Required to specify font color

%\documentclass[a4paper,12pt,openany]{scrbook}

%\usepackage{glossaries}
%\usepackage[xindy]{glossaries}
%\makeglossaries

%\usepackage{odsfile,lmodern}

%\usepackage{luacode}


\usepackage{multicol}

%\usepackage{churchslavonic}
%\usepackage{polyglossia}
%\setotherlanguages{russian,churchslavonic}

\usepackage{mathtext}
\usepackage{makeidx}
%\usepackage{imakeidx}
\usepackage{longtable}
\usepackage[tc]{titlepic}

\usepackage{fontspec}

\usepackage[protrusion=true,expansion,
tracking=true,letterspace=50]{microtype}


%\usepackage{academicons}
%\usepackage{xcolor}


%\SetTracking[ spacing = {25*,166, } ]{ encoding = *, shape = sc }{ 25 }


%\usepackage{soulutf8}
%\usepackage{letterspace}

%\usepackage[toc]{appendix}
\usepackage[toc,titletoc]{appendix}
\renewcommand\appendixtocname{Приложения}
\renewcommand\appendixpagename{Приложения}

%\usepackage{xunicode}
\usepackage{xltxtra}

\usepackage[normalem]{ulem}
\usepackage{adjustbox}

\usepackage{polyglossia}
\setmainlanguage{russian}
\setdefaultlanguage{russian}
\setotherlanguage{ukrainian}
\setotherlanguage{english}
\setotherlanguage{greek}
\setotherlanguage{latin}
\setotherlanguage{polish}


\setmainfont[BoldFont={DejaVuSerif-Bold.ttf},
ItalicFont={DejaVuSerif-Italic.ttf},
 BoldItalicFont={DejaVuSerif-BoldItalic.ttf}]{DejaVuSerif.ttf}

\setromanfont[BoldFont={DejaVuSerif-Bold.ttf},
ItalicFont={DejaVuSerif-Italic.ttf},
 BoldItalicFont={DejaVuSerif-BoldItalic.ttf}]{DejaVuSerif.ttf}

\setsansfont[BoldFont={DejaVuSans-Bold.ttf},
ItalicFont={DejaVuSans-Oblique.ttf},
 BoldItalicFont={DejaVuSans-Oblique.ttf}]
{DejaVuSans.ttf}

\setmonofont[BoldFont={DejaVuSansMono-Bold.ttf},
ItalicFont={DejaVuSansMono-Oblique.ttf},
 BoldItalicFont={DejaVuSansMono-BoldOblique.ttf}]
{DejaVuSansMono.ttf}



\usepackage{enumitem}
\usepackage{indentfirst} 

\usepackage{verse}

\usepackage{graphicx}

%\usepackage[pdfpagelabels,unicode,plainpages=false]{hyperref}

\hypersetup{pdftitle=Предки. Петр Семилетов.}

\usepackage{bookmark} 
\usepackage{relsize}

\usepackage{tocloft,calc}
\usepackage[nottoc,numbib]{tocbibind}


\frenchspacing
\righthyphenmin=2


\makeatletter
\@addtoreset{chapter}{part}
\makeatother  


%\titlepic{\includegraphics[width=0.50\textwidth]{cover/04.jpg}}

\title{ПРЕДКИ\\
\textsmaller[2]{редакция 1.0}}
\author{Петр Семилетов\\ORCID:0009-0002-3901-7785 \orcidlink{0009-0002-3901-7785}}
\date{10/02/2025}



\newcommand*{\plogo}{\fbox{$\mathcal{PL}$}} % Generic publisher logo

% –  –  –  –  –  –  –  –  –  –  –  –  –  –  –  –  –  –  –  –  –  –  –  –  –  –  –  –  –  –  –  –  –  –  –  –  –  –  –  –  –  –  –  – 


\begin{document}


\maketitle

\pagestyle{empty}


\newpage

\pagestyle{plain}

%\hrule

%\tableofcontents

%\input{pre/pre.tex}
%\input{plan-right/plan-right.tex}
%\input{plan-right/plan-right-text.tex}

\begin{abstract}
Не только дохристианские, но и доперуновы верования. Ключом к этой двери послужит киевское урочище Наводничи. В статье изложена подробная привязка давних верований к Наводничам и священному источнику Рожнице.
\end{abstract}


\section{СЕ ПОЧНЕМ}

   Есть в Киеве, неподалеку от Выдубичей, о которых я уже поведал ранее в других статьях цикла про язычество Славян, урочище Наводничи. Его нынче мало кто связывает с язычеством. Тем не менее, это одно из давних священных мест, и оно поможет нам прояснить связь некоторых языческих представлений, ныне почти утраченную. Помните про Розеттский камень? Наводничи и есть такой Розеттский камень. А почему, вы поймете по мере повествования. 

   Если рассматривать бытование верований исходя только из данных Повести временных лет, получается следующая картина.

   Князь Владимир до принятия христианства поставил в Киеве на холму вне двора теремного идолов Перуна, Хорса, Дажьбога, Стрибога, Семарьгла и Мокоши. Идол Волоса уже стоял ранее, и неизвестно, кто и когда его поставил. 

   Значит, Рюриковичи прежде Владимира, будучи язычниками, в Киеве всё же не ставили идолов своего пантеона, во всяком случае в дошедшей летописи про это не сказано. 

Стоял идол Волоса. И до появления князя Владимира на страницах летописи, славянский народ Русь (варяги, пришедшие с Рюриком и его братьями, а также потомки тех варягов), веровал в Перуна и Волоса. А остальное славянское население? Те, кто по мере продвижения Рюриковичей в сторону Киева, попадали под их власть? Словене, Поляне, Древляне, Кривичи, Северы. Они тоже как и Русь веровали в Перуна и Волоса и других богов того пантеона?

Ежели в кого другого, или чтили кого-то другого, других, то продолжилось ли такое почитание после появления новых богов? Судя по всему да, и в итоге, уже после принятия христианства, восточные Славяне стали, в народных верованиях, придерживаться причудливому наслоению из «доперуновой» веры, затем варяжских верований народа Руси, и наконец христианства.
 
   Статью «Скотий бог Волос»\cite{volos} я оборвал на многообещающем отрывке из чудовского списка «Слова Григория Богослова», о дохристианских верованиях Славян:

\begin{quotation}
\noindent се же словени начаша требы класти роду и рожаница преже перуна бога их. а преже того клали упырем. требы и берегыням
\end{quotation}

   Первую половину сего можно толковать однозначно – Славяне, прежде чем класть требы Перуну, клали их роду и рожанице. Вот рожаницы – здесь речь идет не о женщинах, которые рожают, а об особых существах. Класть же требы значит приносить жертвы, или, как говорили на старославянском, жряху. Отсюда и слово жрец – не тот, кто жрёт, но приносящий жертву. А вот жрёт, буквально, уже тот, кому приносят жертву. Слово жертва исходно означало просто пища.

   Итак, согласно этому отрывку, вера в рода и рожаницу предшествовала вере в Перуна. Но есть и вторая часть отрывка:

\begin{quotation}
\noindent а преже того клали упырем. требы и берегыням
\end{quotation}

   Это можно толковать двояко. Вариант первый – упырям и берегиням требы клали еще до рода и рожаницы. Вариант второй – род и рожаница или рожаницы это и есть упыри и берегини.

   Пока не будем разбирать про род и рожаницу, давайте об упырях потолкуем.

\section{УПЫРИ}
    
   У нас это слово вытеснено более западным – вампиром, а ведь фольклорный вампир или упырь сильно отличается от образов, которые навязаны художественной литературой, кино и видеоиграми. 

   Представления об упырях, конечно же, разнились от местности к местности, в рамках моей статьи я не могу охватить всё, поэтому буду усреднять. Упырями называли, по большому счету, колдунов, ведьмачей, живых и мертвых.

 Однако не всех. Как и ведьмы, ведьмачи делятся на родимых или прирожденных, и ученых. Иначе говоря, одни рождаются ведьмачами и ведьмами, а другие учатся этому у первых. Родимые ведьмачи и ведьмы обладают способностью к превращению, миролюбивы, и к их помощи прибегали, чтобы исправить козни ученых ведьм и ведьмачей.

   При жизни упыри мало чем отличались от простых смертных, разве что всячески колдовали. Упырей считали старшими над ведьмами. При полетах на шабаши, бывало, за упыря хваталось много ведьм и так они вместе летели. Любопытно, что по-турецки слово ubyr означает ведьму.

   Как считалось в народе, живой упырь обычно человек крепкий и при том краснолицый, в любом возрасте.

   После смерти упырь лежит в гробу, не разлагается и сохраняет тот же замечательно красный цвет лица. Днем он спит в могиле, а ночью вылезает из-под земли, сидит на кресте, летает по воздуху, гоняется за прохожими, идет в селение и пытается зайти в дома, где сосет у людей, особенно у младенцев, кровь.

   Упырям приписывали разные моры и болезни, поэтому если кого-то подозревали в упырстве, то живого сжигали, а могилу мертвого раскапывали и, обнаружив там нетленное тело, пробивали его осиновым колом.

   Мне кажется, подобные верования в живых мертвецов основаны на следующем. Люди, обладающие особыми способностями, могут впадать в столь глубокий транс, что дыхание едва различимо. Примечательно начало завещания Николая Гоголя:

\begin{quotation}
\noindent I. Завещаю тела моего не погребать до тех пор, пока не покажутся явные признаки разложения. Упоминаю об этом потому, что уже во время самой болезни находили на меня минуты жизненного онемения, сердце и пульс переставали биться… Будучи в жизни своей свидетелем многих печальных событий от нашей неразумной торопливости во всех делах, даже и в таком, как погребение, я возвещаю это здесь в самом начале моего завещания, в надежде, что, может быть, посмертный голос мой напомнит вообще об осмотрительности.
\end{quotation}

   Как известно, окружавшие Гоголя люди, прекрасно знавшие об его завещании, которое было к тому же опубликовано еще прижизненно, похоронили писателя заживо\footnote{В роду Гоголя, его прапрадедом, был полковник Василий Танский, брат знаменитого упыря Антония Михайловича Танского.}.

   Итак, подобные люди, иногда слывшие колдунами, впадали в продолжительный транс. Такой транс, кстати, необходимая составляющая одного из видов оборотничества, когда разум человека подключается к телу животного и управляет им, в то время, как тело лежит. 

   Снорри Стурлусон в «Круге земном», собрании преданий Исландцев, пишет о скандинавском боге Одине. Один ведь был смертным человеком, его обожествили посмертно. И вот что написано про Одина:

\begin{quotation}
\noindent Один мог менять свое обличье. Тогда его тело лежало, как будто он спал или умер, а в
это время он был птицей или зверем, рыбой
или змеей и в одно мгновение переносился в
далекие страны по своим делам или по делам
других людей.
\end{quotation}

   Подобное состояние мы называем трансом. %Славяне в старину использовали для его обозначения слово "мечта". Позже мечта приобрело другое, более отвлеченное, но всё равно близкое значение.
   О блуждании же тела упыря после смерти, то вероятно, очнувшись от транса, колдун находил в себе силы выбраться из могилы, и конечно же, побродив немного по кладбищу, отправлялся к себе домой, где и пугал домочадцев.

   Однако в народе кое-где считали, что телом такого покойника управляет уже не его душа, но вселившийся бес.

   Упырями или одминами, то есмь подменышами, называли также детей, как считали собиратели фольклора, больных рахитом. У них было, как писали, размягчение костей, большие головы, длинные тонкие ноги. Эти дети не могли ходить, зато имели дар ясновидения, и рано умирали. Про них после смерти ничего плохого не говорили, но была молва, что такие упыри сразу после смерти снова рождаются в другом, подобном теле, и опять дают предсказания. Любопытные отголоски прежних верований в христианской среде.

   Если мы вернемся к первым представлениям об упырях, что это колдуны, люди с особыми способностями, то пример Одина показателен. Одина при жизни можно считать колдуном, упырем. Мертвого Одина-упыря обожествили. Что же, может и Славяне обожествляли неких местных, знаменитых упырей, и клали им требы. Что до скандинавских богов-Асов, из числа коих был Один, то они ведь пришли из Азии, потому и Асы. Но это отдельный разговор.

   Вернемся же к чтению и разбору «Слова Григория Богослова».

   Дальнейшие его пояснения должны были, по мнению сочинителя, прояснить дело о дохристианских и даже доперуновых верованиях:

\begin{quotation}
\noindent по святем крещении. перуна извергоша. а по христа бога яшася. но и ныне по украинам молятся ему. проклятому богу перуну хорсу мокоши вилам
\end{quotation}

   После святого крещения, низвергнув Перуна, придя к Христу богу, люди и ныне, по окраинам, молятся ему – проклятому богу Перуну, Хорсу, Мокоши, вилам.

   Запомним упоминание вил! Читаем дальше:

\begin{quotation}
\noindent и то творят отай сего же не могут ся лишити проклятого ставленния.
\end{quotation}

И держатся этой веры тайно, не могут оставить проклятый обычай. 

\begin{quotation}
\noindent вторые трапезы роду и рожаницам на великую прелесть. верным крестьяном. и на хулу святому крещению и на гнев богу 
\end{quotation}

   Ставят вторые трапезы роду и рожаницам, обманывая, вводя этим в заблуждение верных христиан, на хулу святому крещению и на гнев богу. О чем идет речь? Про обычай оставлять кушанья для умерших предков. Обычай сей соблюдался во многих случаях. 

   Под Рождество оставляли на ночь ужин, включавший в себя кутью. Подобный обряд был приурочен ко многим другим праздникам. На поминках ставили для мертвых родственников кушанья под стол. Оставляли еду на могилах во время Гробков, иначе говоря Радуницы. В Радунице мы видим корень – род.

   На христианский праздник Рождество люди едят кутью – а кутья это традиционное блюдо  поминок и похорон.

   Раньше кстати говорили и писали не Рождество, но Рожество. А что сказано в «слове Григория Богослова»? Ставят вторые трапезы роду и рожаницам. Праздник Рождество, Рожество. И рожаницы. Корень-то один. А кто такие эти рожаницы? Это не женщины-роженицы, нет. Это покойные предки или родственники.

   Почему ставить им второе кушанье неправильно с христианской точки зрения?

\section{ПРЕДСТАВЛЕНИЯ О ЗАГРОБНОЙ ЖИЗНИ ПО БИБЛИИ. ДУША И ДУХ}

   Слово «душа» употреблялось в старину в совсем ином значении, нежели сейчас. Ныне под душой подразумевают саму личность, которая может существовать отдельно от тела. Нечто иной природы, вечное и не относящейся к физической составляющей мира.

    Но как быть, например, с подзабытым словом «душегубец»? Если душа бессмертна, то как можно ее погубить? И почему счёт людей, например крепостных крестьян, велся по душам – у помещика столько-то душ?

    Ответ прост – раньше душой называли живое существо, человека, зверя или даже рыбу. Возникает закономерный вопрос – когда в переводе Библии говорится о душе, то в каком смысле? В современном или давнем?

   Необходимо прояснить этот вопрос, чтобы сравнивать с языческими представлениями о жизни за гранью смерти.

   Большая часть того, что люди считают христианскими представлениями о загробной жизни, не исходит из Библии, но составляет надстройки, оставшиеся от языческих верований, а также предлагаемых той или иной церковью. Так, католическая церковь предлагает понятие чистилища, о котором Библия молчит.

   Само по себе слово «душа», вне Библии, исходно означало живое состояние. Есть душа в теле – жив, нет – мертв. Однокоренные слова – дышать, душа, дыхание, дух. Когда человек переставал дышать, люди говорили – из него ушла душа, или – испустил дух, то есть дух выпустился из него, буквально – ушло дыхание. Дыхание – наиболее очевидный признак живого существа, поэтому сей признак люди сочли воплощением личности, разума.

    Итак, по старинным понятиям, душа это живое существо и одновременно жизнь в нем. В латыни для обозначения души используется слово «anima». А животное на латыни будет «animal», то есть производное от «anima».
  
   Сравним латинское «anima» с греческим «άνεμος» (анэмос) – ветер. 

   В славянских переводах Библии слово «душа» в большинстве случаев используется в значении живого человека или другого животного. Например, Бог дает заповедь Моисею:

\begin{quotation}
\noindent всяка душа, яже аще снест кровь, погибнет душа та от людий своих
\end{quotation}

   То есть всякого человека, который съест-выпьет кровь, должны убить свои же соплеменники.

   В Библии именно в смысле что душа это человек или вообще любое живое существо говорится грозное – да истребится душа, да умрет душа, да сотрется душа, потребится душа и так далее.

   В книге Бытия сказано:

\begin{quotation}
\noindent И создал Господь Бог человека из праха земного, и вдунул в лице его дыхание жизни, и стал человек душею живою.
\end{quotation}

   Здесь дыхание сопоставлено с жизнью тела.  Так совпало, что славянское слово «душа» издавна само по себе обозначает обозначает живое существо. В Библии душа, жизнь сопоставляется не только с дыханием, но и с кровью.

«Кровь бо есть душа его» – сказано во Второзаконии 12:23.

    Как и в случае дыхания, было вполне очевидно, что когда живое существо истекало кровью, оно лишалось жизни.

   Душа считалась чем-то, что содержится в живом человеке, делает его живым, а когда она в виде дыхания или крови покидала тело, наступала смерть.

   Помимо равенства между словом душа и живое состояние тела, плоти, слово «душа» иногда употребляется в Библии в значении «личность».

    Если вы внимательно прочитаете Библию, вы не найдете там почти ничего о бессмертии души. Такие представления начинает высказывать только Иисус. Например:

От Матфея 10:27:

\begin{quotation}
\noindent И не убойтеся от убивающих тело, души же не могущих убити: убойтеся же паче могущаго и душу и тело погубити в геенне. 

И не бойтесь убивающих тело, души же не могущих убить; а бойтесь более того, кто может и душу и тело погубить в геенне. \end{quotation}

   Механизм загробной жизни, в Библии, описан скупо, просто и вместе с тем противоречиво. Преимущественно представляется так – если человек умирает, и ежели он не святой праведник, то остается мертвым до второго пришествия Христа, после чего следует воскрешение из мертвых и страшный суд. Святой же праведник воскресает сразу и пребывает рядом с Христом в царстве Божьем. Что происходит с личностью в промежуточном состоянии между смертью и воскрешением, в Библии не сказано. После Страшного суда, одни люди идут в геенну огненную, другие в Царствие Божие.

   Небольшой отступление.

   В Ветхом Завете, после сообщения о том, как кто-то умер, иногда прибавляется, в зависимости от давности перевода: «и приложися к людем своим» или «и приложился к народу своему». Это сходно с известными выражениями – отошел к праотцам, отправить к праотцам. Конец отступления.

   Повторюсь – слово «душа», в большинстве случаев, используется в Библии для обозначения живого существа.

   Также в славянском переводе Библии используется слово «дух» и некие «духи». На латыни в Библии там где «дух» стоит слово «spiritus», то есть буквально переводится как дыхание. А славянское слово «дух» это и дыхание, и запах.

   Библейское понятие «дух» можно толковать как личность, отделенную от тела и способную, входить в разные тела с целью совершения действий либо разговора с окружающими. Это способен выполнять как некий простой дух, так и дух библейского Бога – Иеговы. В последнем случае речь идет о пророчестве – когда дух Бога говорит, прорекает, пророчествует устами человека.

   Итак, дух представляется в Библии как некий разум, личность, состоящую из дыхания, ветра, получается – газообразную. Возможно, библейский Дух Божий, на каком-то этапе представлялся людям именно дующим в природе ветром – он был незрим, но находился всюду, люди ощущали его непосредственно. Люди видели облака в небе. Люди попадали в бури на море и на суше.

   Со временем это буквальное понимание духа как носителя жизни в теле изменилось. Богословы отняли у понятий души и духа их материальную природу и заменили ее ничего не значащим обозначением «духовный». Это всё равно что если попросить человека объяснить, что такое красный цвет, а вы получите ответ, что красный это красный. Так и здесь. У духовного было буквальное значение – можно сказать, газообразный, ибо дыхание, ветер. Богословы стёрли буквальное значение, предложив новое – духовный это духовный.
   
\section{ПРЕДСТАВЛЕНИЯ НАРОДА О ДУШЕ}

   Мы ознакомились с представлениями о душе, выраженным в самой Библии. Что же думал о душе народ, которому принесли христианство? Принесли не вместо язычества, а в добавление к оному?

    Проводники христианства, сами толкующие несомое учение на разные лады, внесли в представления народа о мироздании невероятную сумятицу. Судя по всему, до христианства народ верил, что после смерти человек, оставив прежнее своё тело, продолжает существовать как бы в другом режиме, иногда сохраняя пространственно-временную связь с живущими.

    Но в христианское время такое могло считаться, подчеркиваю, считаться справедливым только для умерших без крещения, то бишь младенцев, а также для умерших неестественной смертью – например утопленниц, утопленников, самоубийц и так далее. Кое-где считалось, что у человека человек должен прожить определенный срок, и если он умирает раньше, то как бы застревает в ином состоянии на остаток непрожитого срока.

   Для умерших неестественной смертью было много названий. У западных народов, например, считалось, что такие люди могут стать эльфами. А у Славян – русалками.

Это сейчас в сознание внедрили образ, что русалка – девушка с рыбьим хвостом. Произошло сие от нерадивости переводчиков, которые, в самом деле рыбьехвостых морских дев из западных сказок переводили как русалок.

Представления же о русалках славянского народа нельзя толком свести воедино, слишком разнится их образ от места к месту. То они маленькие, как дети, в красных рубашках, бегают по полю. То волосатые, качаются на ветках, имеют, как описывали крестьяне, «железные цыцки». То плавают в воде и выглядят как обычные женщины и мужчины. То – «все белые».

   Однако несмотря на внешние различия, суть у русалок одна – ими считались умершие, продолжавшие в ином виде жить люди, женского и мужского пола. Для покойных некрещеных младенцев было еще одно название – потерчата, а вероятно также и упомянутые в разбираемой нами рукописи берегини.

   Вообще именований русалок множество – криницы (от того, что сидят на криницах – колодцах), богини, семилетки – от того, что подобны семилетним детям. Есть связанный с русалками праздник – Семик, четверг седьмой недели после пасхи. В некоторых местах Семик называли просто Русалка. На Семик поминали покойников, умерших неестественной смертью или некрещеных детей, и вообще ходили пировать на кладбище, как на Радуницу, иначе говоря Гробки.

   Известна также особая неделя – ее называют Зеленой, Клечальной, Русальной, Троицкой – ибо начинается за днем Троицы. Говорят еще – Гряная неделя. Праздник Троицы привязан к дате Великодня, Пасхи – на 50 день после нее. Поэтому расположение русальной недели в году как бы плавает. 

Гряная неделя считается днями, когда русалки много взаимодействуют с людьми. На этой неделе были отдельные праздники – на ней справляли проводы и похороны русалок, Костромы, Родительский день, Русалчин Великдень, он же Навская троица. Можно считать русальные дни – поминками, но поминками веселыми, с песнями и плясками.

   Историк С. М. Соловьев в 1849 году задал любопытный вопрос:

\begin{quotation}
\noindent Народ теперь верит, что русалки суть души младенцев, умерших без крещения. Но когда все славяне умирали некрещеными, то души их все должны были становиться русалками?
\end{quotation}

В самом деле, не таково ли было дохристианское представление о том, что может случиться за гранью смерти? Далее Соловьев продолжает:

\begin{quotation}
\noindent Русальные игры суть игры в честь мертвых. Кроме русалок, души умерших, мертвецы, были известны еще под именем Навья.
\end{quotation}

Был даже праздник Навский великдень – совмещенный с чистым Четвергом, первым Четвергом после Пасхи. В Навский, или Мертвецкий Великдень (четверг до или после Великодня, Пасхи) люди готовили поминальный обед по своим покойникам, и считалось, что это один из дней, когда мертвые возвращаются на землю. В некоторых местностях поверье это было связано не просто с мертвыми, но с русалками. Ходило множество быличек о том, как ночью в этот день людей преследовали ожившие мертвецы.

Не надо путать Навский великдень с Русальчиным великоднем, то четверг на Троицкой – Гряной неделе.

   Слова «нави» очень древнее, к 19 веку уже никто так мертвецов не называл, оно осталось в народе только для обозначения праздника, да в некоторых областях Украины русалок именовали нявками или мавками – очевидно производными от навей.

   Нави встречаются однако в стародавних письменных источниках.

   Радзивилловская летопись за лето 6598 от библейского сотворения мира сообщает:

\begin{quotation}
\noindent приведе Янка митрополита Иоана скопьчину егоже видевше людье рекоша се навье пришел год пребыв оумре
\end{quotation}

Люди видели скопца Иоана и говорили – это навье пришел. Год пробыв, он таки умер. А в Ипатьевской летописи вместо «навье» стоит просто «мертвец». 

В сочинениях Иоанна Ексарха Болгарского находим:

\begin{quotation}
\noindent Яко и навь из гроба исходящи.
\end{quotation}

    Итак, запомним – навь, навье это покойники, однако слово это кажется всегда применялось по отношении к тем мертвецам, которые проявляют себя живыми, действуют.

   Снова повторю выдержку из поучения:

\begin{quotation}
\noindent се же словени начаша требы класти роду и рожаница преже перуна бога их. а преже того клали упырем. требы и берегыням
\end{quotation}

   Славяне приносили жертвы роду и рожаницам, прежде того, как стали делать это Перуну, а прежде клали требы упырям и берегиням. 

   Попробуем выяснить, кто такие эти род и рожаницы, и как толковать последнюю часть. Упыри и берегини сопоставлены с родом и рожаницами? Или же сочинитель имел в виду, что им требы клали еще до того, как стали класть роду и рожаницам?

   Ранее я намеренно, чтобы не запутывать дело, сокращал выдержку из Слова Григория Богослова по Чудовскому списку. Полностью же читается так:

\begin{quotation}
\noindent се же словени начаша требы класти роду и рожаница преже перуна бога их. а преже того клали упырем. требы и берегыням по святем крещении. перуна извергоша. а по христа бога яшася. но и ныне по украинам молятся ему. проклятому богу перуну хорсу мокоши вилам и то творят отай сего же не могут ся лишити проклятого ставленния. вторые трапезы роду и рожаницам на великую прелесть. верным крестьяном. и на хулу святому крещению и на гнев богу веруют упырем и младенци знаменают мертви и берегиням их же нарицают 7 (9 по другому списку) сестрениц. а друзии веруют в сварожитьца. и артемида и артемидию имже члвеци невегласи молятся. и куры им режут и то блудивше тоже сами ядять.
\end{quotation}

Выделим отсюда новый, весьма мутный отрывок:

\begin{quotation}
\noindent веруют упырем и младенци знаменают мертви и берегиням их же нарицают 7 (9 по другому списку) сестрениц
\end{quotation}

Попробуем разобрать. Веруют в упырей – об этом мы уже поговорили. 

«И младенци знаменают мертви». Что значит знаменают? Однокоренное слово – знамение. Откроем словарь Даля:

\begin{quotation}
\noindent ЗНАМЕНАТЬ церк. знаменовать что,  означать знаком, налагать знак, метить, помечать, обозначать; подавать знак, маячить;  означать что собою, изображать, представлять, выражать, заключать в себе значение,
смысл;  предвещать.   
\end{quotation}

И приводит пример:

\begin{quotation}
\noindent Багровый закат знаменует ветреный день.
\end{quotation}

То есть можно подумать, что мертвые младенцы что-то предвещают, или же напротив, речь идет о гадании на мертвых младенцах. Звучит дико, однако вспомним, что писал Лев Диакон, византийский историк, живший в 10 веке, в своей книге «История». Кроме прочего он говорит там о соотечественниках нашего князя Святослава Игоревича, именуя его Сфендославом, а собственно его соотечественников Росами и Скифами. 

Прежде чем зачитать отрывок, отмечу, что Греки все имена переиначивали на свой лад, поэтому вполне вероятно, что даже в известных нам греческих мифах действовали Славяне, коих Греки наделили греческими или огреченными именами.

   Итак, отрывок. Напоминаю – Скифы это Росы, воины князя Святослава. Наука считает, что Скифы – некие ираноязычные племена, но вот современник скифов Лев Диакон вполне определенно именует скифами Росов.

   Отрывок:

\begin{quotation}
\noindent И вот, когда наступила ночь и засиял полный круг луны, скифы вышли на равнину и начали подбирать своих мертвецов. Они нагромоздили их перед стеной, разложили много костров и сожгли, заколов при этом по обычаю предков множество пленных, мужчин и женщин.

Совершив эту кровавую жертву, они задушили грудных младенцев и петухов, топя их в водах Истра. 

Говорят, что скифы почитают таинства эллинов, приносят по языческому обряду жертвы и совершают возлияния по умершим, научившись этому то ли у своих философов Анахарсиса и Замолксиса, то ли у соратников Ахилла. Ведь Арриан пишет в своем «Описании морского берега», что сын Пелея Ахилл был скифом и происходил из городка под названием Мирмикион, лежащего у Меотидского озера. \end{quotation}

Мэотидское озеро – это Азовское море. После сражения, воины князя Святослава – Росы, они же Скифы, собрали своих убитых на поле боя и принялись сжигать их на кострах, попутно, по обычаю предков, убивая пленных.

Совершив эту кровавую жертву, – продолжает Лев Диакон, – «они задушили грудных младенцев и петухов, топя их в водах Истра». 

Истр это Дунай.

Налицо чудовищный обряд, при котором умерщвляют младенцев и петухов – возможно душат и затем топят в реке, или просто топят. Быть может, именно это согласуется со словами из слова Григория Богослова «и младенци знаменают мертви»? Тем более что далее он тоже говорит про умерщвление кур:

\begin{quotation}
\noindent от убогая курята, яже на жертву идолом режются, инии в водах потапляемы суть\end{quotation}

   А может быть, Лев Диакон всё выдумал, чтобы очернить язычников? Да нет. Люди были и остаются жесточайшими существами на планете, разница только в законах, по которым они живут.

   Про обрядовое убийство младенцев Славянами писали также историки средневековья Псевдо-Кесарий, Герборд и Эккерхард. Впрочем, сие практиковали и другие народы и религиозные группы. Кроме того, вероятно, не все народы Славян совершали такие ритуалы. Мало данных для суждения.

   Что же, одно толкование выражения «и младенци знаменают мертви» мы приблизительно получили. Но ведь есть еще слово «знаменито», означает оно «известно», вернее «прославленно». Тогда выражение приобретает иной смысл – прославляют мертвых младенцев.

   Язычник князь Святослав и его дружина связаны с известным пантеоном поименованных богов – Перуна, Волоса. Но мы знаем из того же поучения, что 

\begin{quotation}
\noindent се же словени начаша требы класти роду и рожаница преже перуна бога их. а преже того клали упырем. требы и берегыням
\end{quotation}

Вот эти род, рожаницы, упыри, берегини и мертвые младенцы идут как даже не дохристианские, а доперуновы верования тех Славян, которые не разделяли почитание пантеона Рюриковичей.

   Повторение того же утверждения в Слове Григоря Богослова, напомню, таково:

\begin{quotation}
\noindent веруют упырем и младенци знаменают мертви и берегиням их же нарицают 7 (9 по другому списку) сестрениц
\end{quotation}

Тут неясно, то ли младенцы мертви это и есть берегини, или же отдельно младенцы, и отдельно берегини. И берегинь называют еще 7 сестрениц – то есть берегини не просто бесчисленны, но их семеро или 9 сестер. Подобное представление есть о вилах в «Слове некоего христолюбца и ревнителя по правой вере»:

\begin{quotation}
\noindent и в волы (вилы) их же числом 39 сестрениц глаголють бо невигласни то все мнять богынями
\end{quotation}

Вилы были известны южным и западным славянам, как девушки, что умерли некрещеными или не своей смертью – словом, соответствовали русалкам. У хорватов вилы являлись к новорожденным детям и предрекали будущую судьбу. У Ф. Буслаева есть статья «О сродстве славянских вил, русалок и полудниц с немецкими эльфами и валькириями», где он отмечает:

\begin{quotation}
\noindent Как вилы и русалки почитаются душами некрещеных младенцев, так и Эльфы не только души умерших – потому-то они и пляшут на кладбище и празднуют смерть человека – но и падшие с неба ангелы
\end{quotation}

Мы видим тут весьма точную формулировку народных представлений об этих существах. Церковь говорила о них как о бесах, а народ помнил старые свои верования, что это люди, которые умерли и теперь живут в каком-то ином режиме существования.
 
   Может, в веке 19-м народ помнил и о берегинях? Кое-где по Украине девушки изображали Перегиню – поднимали над головой руки, держа в них клубок, как бы голову, и всё это покрывалось одеждой – получалось высокое существо с вытянутой шеей.

Какие синонимы собрал к берегиням Даль в своем словаре?

\begin{quotation}
\noindent Берегиня  стар.  берегуша, лопаста, шутовка, водяная, русалка.
\end{quotation}

    Подытожим и двинемся дальше. Род и рожаница, берегиня, русалка – одного поля ягоды, умершие люди, но действующие, думающие, само их существование противоречит библейской модели мироздания, но в них верили, их встречали, их почитали.

   Теперь я расскажу вам о забытом ныне славянском языческом святилище, именуемом Рожница.

   Как я на него вышел? 

   В девяностых я много гулял со своей собакой, Бобиком, с Бастионной вниз по Зверинецкому холму в парк Примакова. У подножия горы, склон которой опоясывает улица Мичурина (Ломаковская) был родник – один из, остальные связаны дренажкой. Родник столь знаменитый, что по крайней мере в 1980-е и девяностые люди со всего Киева отправлялись туда за водой с бидонами и банками, к такой трубе, откуда истекал ощутимый поток студёной воды. То есть шла об источнике народная молва и не прекращалась.

   Потом конечно там соорудили автозаправку, но есть сведения что родник поныне сохранился, за оградой этой автостоянки.

\begin{center}
\includegraphics[width=\linewidth]{pix/IMG_20150411_131729.jpg}
\end{center}

   Где это? А вот покажу вам на карте.

\begin{center}
\includegraphics[width=\linewidth]{pix/map01.jpg}
\end{center}

   Чуть дальше от перекрестка бульвара Дружбы Народов, Старонаводницкой и Лаврской улиц. С одной стороны Зверинецкий холм, с другой холм с огромным памятником-статуей Родины-матери.  Между ними удолье – низовья урочища Наводничи, или же овраг Наводницкая балка. Конец ее у моста Патона, а начало наверху  улицы Старонаводницкой. Вдоль оврага, в подземном коллекторе протекает Наводницкий ручей, раньше он бежал по поверхности, и где-то на нем было озеро Проклятое. Также, в 17 веке, овраг называли Душегубицей.

    Когда я увлекся краеведением, то стал частенько разбирать давние земельные документы. В 16-17 веках в земельных грамотах про монастырские владения ходило слово Наводничи или Неводничи, относительно этих мест.

   И один из земельных споров между Лаврой и Выдубицким монастырем был завершен 1 мая 1613 года подписанием соглашения, а произошло это

\begin{quotation}
\noindent на месцу Рожнице там, где делится долина Зверинская Неводничи, з озерка поточком з ровчаком в Затон идучим.
\end{quotation} 

   Любопытно, что духовные лица подписали    подписали мировую не допустим в ратуше, а на природе, в некоем месте, урочище под названием Рожница. Почему? Что такого особенного в этом месте? Может, была устойчивая традиция заключать там договоры? Прежде заключались договоры в священных урочищах, дабы сверхъестественное существо было свидетелем договора и покарало нарушителя. Но ведь подписывающие стороны – представители христианской церкви. Загадка.

   И указано, что Рожница находится там, где Наводницкая долина разделяется. 

   Речь идет скорее всего именно о месте родника, где ныне автозаправка\footnote{50°25'25.7"N 30°33'46.9"E}. Или ближе к перекрестку. Сейчас открытых родников нет, к тому же славой был овеян именно описанный мною родник, хотя конечно возможно, что какой-то другой родник неподалеку сгинул, и внимание народа переместилось на близлежащий.

   В любом случае ничего иного примечательного там не прослеживается ни на старинных картах, ни в документах. Поэтому с большой степенью вероятности предположим, что Рожницей называли место, где этот родник. 

   Рожница. Что означает название?

   Две версии. 

   Первая – название произошло от слова рог, то бишь угол. Рожница, на углу, на розе. И с большой натяжкой можно сказать, что родник находится на углу.

   Но что если Рожница происходит именно от рожаниц?

   Давайте посмотрим, как всё любопытно складывается.

    Рожаницы, род. Вот у нас над урочищем-родником Рожницей нависает с противоположной горы статуя Родины-матери. Родина... Рожаница... Рожаницы связаны с культом предков. Но в постаменте статуи полнится эхом зал, исписанный именами тысяч погибших в войне, это же создавалось как музей Великой Отечественной войны.

   Рожница – а рядом такое говорящее название статуи, и такая тесная связь с почитанием мертвых. Совпадение?

   А в каком урочище находится Рожница? 

   В Наводничах. Навье место, не правда ли?

   Еще одно совпадение? 

   Рожница – Родина мать.

   Навы – Наводничи.

   Осталось русалок поблизости найти. Но, по сути русалки и есть навы. И рожаницы.

   Если мы найдем привязку к окрестностям еще и русалок, то это будет уже не совпадением, а подтверждением закономерности.

   Что же.

   Давно, над Наводничами, в 360 саженях от Васильковской башни Печерской крепости, стояла Пятницкая башня с одноименными воротами, близ которых была церковь святой Параскевы. От этих ворот, по росписи Киева за 1682 год, ездили на Зверинец, Наводницкий перевоз и Выдубицкий монастырь.

   Святая Параскева была известна в народных поверьях как святая Пятница. Греческий день недели «параскэвы» (Παρασκευή) соответствует нашей пятнице. А в «Описании Лавры» 1837 года находим указание, что до устроения Печерской крепости в 1706 году, около Лавры было 955 крестьянских дворов, относящихся к приходам Спасскому, Феодосиевскому и прочим, в том числе 155 дворов в Пятницком приходе – стало быть, название Пятницы закрепилось за церковью вместо Параскевы, а церковь появилась раньше крепости.

   В фольклоре, Пятница была ходячим олицетворением запрета работы в упомянутый день. Нарушение запрета – если люди пряли, расчесывались, трусили сажу, квасили, стирали – Пятницей жестоко каралось.

   Рассказывали о встречах с нею. Голая, с распущенными волосами, догоняла людей. Праздник святой Пятницы отмечали обычно у воды, в пятницу на 9-10 неделю после Пасхи, что близко к Зеленой неделе, столь известной активностью русалок. Кстати, в церквях в шкафчиках ставили деревянные раскрашенные статуи Пятницы, которым молились.

   Зеленин в книге «Очерки русской мифологии: Умершие неестественной смертью и русалки»\cite{zeleninrusalki} пользуется в качестве одного из источников рукописью И. Голяховского из города Валок, Харьковской губернии, 1870 года. Голяховский писал, что Пятницу многие местные причисляют к русалкам. Существовал у восточных славян и обряд, прямо связанный с, как называли, «русалкой Пятницей». Пятницу считали русалкой!

   Так что же это? В урочище Наводничи, чье название связано с навками-русалками, стояла церковь в честь русалки Пятницы, и было урочище Рожница, которая, если Рожница это Рожаница – тоже нечто вроде русалки. Всё это вместе уже определенно указывает на священное для язычников значение Наводничей и связь этой местности с русалками, и основной священный источник в ней – Рожница, где спустя много лет после введения христианства, даже представители христианских монастырей подписывали договоры!

\section{ПОСЛЕСЛОВИЕ}

Судя по всему, вообще вся та часть правого берега Днепра была богата на языческие святилища. Южнее Наводничей – Выдубичи. В других статьях цикла («Как Владимир идолов низвергал»\cite{kak} и «Скотий бог Волос»\cite{volos}) я рассказывал о том, что Выдубичи были урочищем, где ежегодно происходил завершающий этап давнего ритуала, сходного с проводами русалок. Выше по течению от Наводничей – сейчас сверкает золотыми куполами Лавра, но что было там прежде? По крайней мере мы знаем о Варяжских пещерах, в которых монахи встречали неких бесов, что снова отсылает нас к языческим верованиям, поскольку для христианства всё языческое это бесовское.

   Были в Киеве и другие подобные места, причем на левом берегу Днепра. Но об этом – в другой статье.


\newpage

%\backmatter

%\appendix
%\addappheadtotoc


%\input{appendix/dict.tex}

\begin{thebibliography}{999}
\bibitem{zeleninrusalki}
\emph{Избранные труды. Очерки русской мифологии: Умершие неестественною смертью и русалки}. Зеленин Д. К. Индрик, Москва, 1995.

\bibitem{kak}
\emph{Как Владимир идолов низвергал}. Петр Семилетов, DOI 10.5281/zenodo.14212537

\bibitem{volos}
\emph{Скотий бог Волос}. Петр Семилетов, DOI 10.5281/zenodo.14563494
\end{thebibliography}


\end{document} 